\documentclass[11pt,a4paper]{article}
\usepackage[margin=2.3cm]{geometry}
\usepackage[T1]{fontenc}
\usepackage{mathptmx}
\usepackage{enumitem}
\usepackage{hyperref}
\setlength{\parskip}{4pt}\setlist{leftmargin=*,itemsep=2pt}
\newcommand{\rc}[2]{\par\medskip\noindent\textbf{#1}\par\noindent\textit{Response: Done.} #2\par}
\begin{document}
\begin{center}\Large\textbf{Response to Reviewers}\end{center}
\noindent Manuscript NEURI-D-26-00299, \emph{Neuroscience Informatics}: ``An Explainable Stacking Framework for Multiclass Sleep-Disorder Risk Stratification''.

\medskip\noindent We thank the Editor and the five reviewers for their careful and constructive reports. The original manuscript has been revised in place; its structure, figures and layout are retained. Every analysis requested was carried out, and every model-dependent table, figure and in-text number was regenerated from a single released pipeline, so that the whole manuscript is internally consistent. Section and table names below refer to the revised manuscript.

\section*{Overview of the revision}
\begin{itemize}
\item \textbf{Verified data and outcome coding.} The source file was audited directly: 374 records, 219 literal ``None'' labels, 77 insomnia, 78 sleep apnea, no empty outcome fields. The audit also quantified repetition: 242 exact duplicates (132 unique rows, 109 distinct predictor vectors) and 22 predictor vectors with conflicting labels, which bounds attainable accuracy at 93.3\%.
\item \textbf{One primary model throughout.} The two-learner RF--XGBoost stack with a multinomial logistic meta-learner is the model used in every analysis: validation, calibration, ablation, repeated splits, SHAP and LIME.
\item \textbf{Rigorous validation.} 10$\times$5 repeated stratified-grouped cross-validation (grouped by predictor vector), repeated nested cross-validation with tuning inside training partitions, a unique-row analysis, cluster-bootstrap confidence intervals, Brier score, log-loss and calibration curves, and paired tests (exact McNemar, Nadeau--Bengio corrected $t$-test, paired cluster bootstrap).
\item \textbf{Broad baselines.} Twelve comparators under identical folds: majority class, depth-3 decision tree, logistic regression, RBF SVM, MLP, FT-Transformer, random forest, extra trees, HistGradientBoosting, LightGBM, CatBoost and XGBoost, plus RF--XGBoost soft voting.
\item \textbf{Audited explanations.} Ten independent permutation-SHAP runs, block and collinearity analyses (Spearman, VIF), LIME with identifiers, probabilities and surrogate fidelity, and a bootstrap uncertainty analysis of the ALE surface.
\item \textbf{Reproducibility package.} Code, software versions, seeds, split identifiers, observation-level labels, predictions and probabilities for every model, repeat and fold, and a script that regenerates every table, figure and in-text number.
\item \textbf{Presentation.} All figures redrawn in a uniform publication style, all tables regenerated, equations and typography checked, new references added.
\end{itemize}

\section*{Reviewer \#1}
\rc{1. Verify outcome coding.}{The source CSV was opened and the outcome column read with default missing-value conversion disabled. It contains exactly three strings: \texttt{None} (219), \texttt{Sleep Apnea} (78) and \texttt{Insomnia} (77); no field is empty. ``None'' is therefore a literal category label, which default parsers silently convert to a missing value. The exact import and recoding code is released and described in Section ``Data source and outcome definition'' and Algorithm~1, and the analysis was repeated end to end with the verified coding. Because the label is dataset-supplied, the manuscript treats it as a category of the benchmark and states this in Data and Limitations.}
\rc{2. Reconcile model configurations.}{The ablation and the repeated-split analysis were re-run on the primary two-learner RF--XGBoost stack and fully specified. Section ``Stacking Ablation Study'' evaluates seven configurations (full stack, without RF, without XGBoost, without pass-through, soft voting, RF alone, XGBoost alone) on identical grouped folds; all lie within 0.3 percentage points (89.6--89.9\%). The repeated-split figure now covers 13 seeds (50--650) of the same stack, with exact Clopper--Pearson bands and a numerical results file. All SVM-based analyses were removed.}
\rc{3. Establish reproducibility.}{The release contains executable code, software versions, seeds, split identifiers and a one-command regeneration script. Preprocessing (imputation, scaling, one-hot encoding) is fitted inside each stacking fold on that fold's training part only, through a pipeline object that is cloned per fold; out-of-fold probabilities are built from five inner folds; the test partition is excluded from all selection, and hyperparameters are fixed a priori or tuned in inner folds of the training partition only (repeated nested CV). See ``Experimental protocol'' and ``Reproducibility''.}
\rc{4. Strengthen validation and comparisons.}{Section ``Experimental protocol'' and ``Comparison with ensemble baselines'' report 10$\times$5 repeated grouped CV (50 folds) with cluster-bootstrap 95\% intervals. Stacking reaches 89.6\% accuracy (95\% CI 86.1--92.7; macro-F1 0.87) versus 89.8\% for soft voting, RF and XGBoost, 89.7\% for logistic regression, 88.1\% for a depth-3 tree and 58.6\% for the majority class. Repeated nested CV (inner 3-fold selection) gives 90.0\% for the stack with the same ordering. Fold-to-fold SD is 3.6--4.0 points and repeat-to-repeat SD about 0.5 points. Paired differences with intervals, exact McNemar and corrected $t$-tests are tabulated; the unique-row analysis gives 73.0\% for the stack.}
\rc{5. Verify statistical results.}{All observation-level labels, predictions and probabilities are released, and the paired bootstrap and McNemar test are recomputed from them. On the 75-observation hold-out, the stack and soft voting each classify 72 records correctly (96.00\%; exact 95\% CI 88.8--99.2); the correctness table is $[[71,1],[1,2]]$, McNemar $p=1.000$, and the paired bootstrap interval is $-4.00$ to $+4.00$ points. The text consistently reflects equivalence: the minimum detectable difference for 75 records is about 5.3 points and about 1{,}180 records would be needed to resolve 1.33 points.}
\rc{6. Improve explanation reporting.}{Local explanations identify the observation (Person ID 244, reference insomnia), report predicted probabilities (stack 0.90 / 0.08 / 0.02; voting 0.97 / 0.03 / 0.00) and surrogate fidelity; a table reports Person ID, reference label, probabilities, $R^2$ and seed-to-seed rank correlation for every explained record. Sampling is specified (25 training background records, 35 explanation records) and SHAP stability is assessed over ten independent runs (rank correlation 0.97--0.98, minimum 0.94). The repeated-split bands are exact 95\% Clopper--Pearson intervals, and the underlying numbers are in the released \texttt{seed\_sweep.csv}.}
\rc{Methodological contribution.}{The Introduction and Discussion state the contribution precisely: a leakage-controlled implementation of established components with a quantified evaluation: data audit (duplicates, label conflicts, ceiling), grouped and nested validation, calibration, paired inference, and fidelity and stability audits of SHAP and LIME. The wording ``proposed framework'' was replaced accordingly.}

\section*{Reviewer \#2}
\rc{1. Limited novelty.}{We agree that stacking is established. The manuscript now presents the work as a rigorous, reproducible evaluation methodology for small tabular clinical benchmarks, with the contributions listed in the Introduction and revisited in the Discussion.}
\rc{2. Clinical interpretation limited by data and outcome definition.}{The task is described as classification of a recorded outcome category, with the audit of provenance, duplicates and label conflicts in the Data section, a dedicated Discussion paragraph separating benchmark performance from clinical screening, and the Limitations section. The title, abstract and conclusion use ``computational benchmark'' wording.}
\rc{3. High performance on a small hold-out requires more rigorous validation.}{The hold-out is reported with exact intervals and its leakage structure (68 of 75 test records repeat a training predictor vector), and the primary evidence is the repeated grouped and nested cross-validation described above, with a unique-row analysis (73.0\%).}
\rc{4. Explanation methods do not demonstrate mechanisms.}{Sections on SHAP, LIME, PCA and ALE state that they describe the fitted function, report stability (ten runs), collinearity (VIF up to 38, Spearman up to 0.98), block permutation importance, LIME fidelity and ALE bootstrap uncertainty, and the Discussion states that they do not support mechanistic or causal conclusions.}
\rc{5. Comparison with recent and competitive tabular approaches.}{The comparison now includes an FT-Transformer (feature-token transformer), CatBoost, LightGBM, HistGradientBoosting, extra trees, an RBF SVM and an MLP under identical folds, with paired differences. All reach 88.7--90.0\% grouped-CV accuracy, consistent with the literature on tree ensembles for small tabular data. An evaluation script for the pretrained TabPFN model under the same folds is included in the release, and TabPFN is listed in the Discussion as a natural extension.}
\rc{6. Presentation, formatting, figures, references.}{All figures were redrawn in a consistent publication style (matching fonts, panel labels, colour palette, error bars with defined intervals); tables and in-text numbers are generated automatically from the results files; equations, cross-references and typography were checked, and the manuscript compiles without errors or unresolved references. Relevant references were added, including the comparator and calibration literature.}

\section*{Reviewer \#3}
\rc{1. Separate benchmark outcomes from clinical screening.}{A dedicated Discussion paragraph explicitly separates computational benchmark outcomes from real-world screening requirements (adjudicated labels, prevalence, external cohorts, decision thresholds).}
\rc{2. Expand interpretation of McNemar/bootstrap.}{The Paired comparison section explains that with 75 test records one record equals 1.33 points, the minimum detectable difference is about 5.3 points, about 1{,}180 records would be required to resolve a one-record difference, and gives the correctness table, McNemar test and bootstrap interval. The repeated-CV paired tests (difference $-0.2$ pp, 95\% CI $-1.8$ to $+1.1$) are added.}
\rc{3. Collinearity and attribution differences between architectures.}{The text and a stability table document how near-collinear predictors (systolic/diastolic pressure, Spearman 0.98, VIF 35--38; sleep duration, quality and stress, $|\rho|$ 0.81--0.91) redistribute attribution: the stack is led by diastolic pressure and the voting model by systolic pressure, while the blood-pressure block carries comparable importance in both, and the SHAP--LIME rank correlation of the stack is 0.91.}
\rc{4. Repeated-split analysis used a different model.}{Re-run on the primary RF--XGBoost stack, and stated in the section opening and the figure caption.}
\rc{5. Labelling note in limitations.}{Added to Limitations: ``None'' is a dataset-supplied category, outcomes are not clinically adjudicated, and the possible classification bias is discussed.}
\rc{6. Local explanation reporting.}{The Discussion lists the information to report with every local explanation (record identifier, reference label, predicted probabilities, surrogate fidelity, sampling and seed), and the manuscript's own figures and table follow it.}
\rc{7. Future work.}{Future work lists repeated cross-validation with nested tuning, probability-calibration assessment, and external clinician-adjudicated cohort validation; the first two are already implemented in this revision.}
\rc{8. Suggested references.}{The three suggested works (Cao et al.; Teng et al.; Yin et al.) were added to Related Work and the reference list.}

\section*{Reviewer \#4}
\rc{Major 1. Provenance, synthetic data, ``None'' coding.}{The source file was opened and the coding verified (see Reviewer 1, point 1). The Data section states that the file has no documented collection process and that its repetition structure (242 duplicates, 22 conflicting vectors, 93.3\% ceiling) is consistent with synthetic or templated records; the results are presented as benchmark behaviour.}
\rc{Major 2. Verified results and deposited code.}{All results were generated and checked by the released code; ``supplied'' matrices were replaced by recomputed ones. Code, split indices, observation-level predictions and a regeneration script are deposited.}
\rc{Major 3. Duplicates and tuning.}{Duplicates are audited and handled by grouping all records with an identical predictor vector in the same fold, so no vector appears on both sides of any split. A unique-row analysis is added. Hyperparameters are specified in a table; repeated nested CV selects them in inner folds of the training partition only, and the test partition never informs them.}
\rc{Major 4. Repeated nested CV, CIs, calibration, baselines.}{Implemented: repeated nested CV, cluster-bootstrap intervals, Brier score, log-loss, ECE and a calibration plot, logistic regression and a depth-3 tree, and ten further baselines. Brier scores range from 0.172 (SVM) to 0.209 (tree), 0.186 for the stack, and the hold-out interval (88.8--99.2\%) is reported next to the repeated-CV estimate (89.6\%).}
\rc{Major 5. Align explanation analyses; SHAP stability.}{All explanation analyses now use the primary stack (stack SHAP, LIME, beeswarm, radar) with the soft-voting model retained as a labelled comparison. The radar panel states its model and sample (35 explanation rows). SHAP stability was assessed over ten independent runs, collinearity quantified, and fidelity reported.}
\rc{Minor 1. ``Risk stratification'' and circularity.}{The task is described as classification of a recorded status; the overlap of sleep duration and quality with the insomnia definition is discussed under label circularity, and a predictor ablation removing sleep duration, quality and stress shows accuracy essentially unchanged.}
\rc{Minor 2. Novelty/``proposed framework''.}{Reworded as above.}
\rc{Minor 3. Table of ANN accuracies.}{The ANN comparison is presented as descriptive literature context, with the matched comparison made through the grouped-CV table with intervals.}
\rc{Minor 4. Observation index, probability, bands, train/test closeness.}{The LIME figure carries Person ID 244 and probabilities; bands are exact Clopper--Pearson intervals. Training/test closeness (0.928 / 0.925) is explained in the text: on average 88.5\% of test predictor vectors also occur in the training partition.}
\rc{Minor 5. ALE grid.}{The surface was quantified with 200 bootstrap resamples (mean 95\% band half-width 0.059 against a range of 0.204); the effective grid (6 $\times$ 4) and cell occupancy (eight cells hold 91.3\% of records) are reported, and the panel is described as descriptive.}
\rc{Minor 6. Condense Related Work and repeated caveats.}{The Related Work table was relabelled as literature context and the ANN paragraph shortened; the clinical-validity caveat is stated once in each of the Introduction, Discussion, Limitations and Conclusion.}

\section*{Reviewer \#5}
\rc{1. 219 blank outcomes.}{Verified: they are the literal label ``None'' (see Reviewer 1, point 1); import, cleaning and recoding code is released.}
\rc{2. 97.33\% on 75 observations.}{The hold-out was re-executed with the released pipeline and is reported as 72/75 (96.00\%, exact 95\% CI 88.8--99.2; macro-F1 0.94) together with per-class results. Interpretation rests on 10$\times$5 grouped CV (89.6\%, 95\% CI 86.1--92.7) and nested CV (90.0\%), which are stable across repeats (SD 0.5 points).}
\rc{3. Code, versions, seeds.}{Complete analysis code, package versions, seeds, preprocessing and model-fitting scripts are in the release.}
\rc{4. Consistent model in ablation and repeated analysis.}{Both analyses use the primary two-learner RF--XGBoost stack (see Reviewer 1, point 2).}
\end{document}
